\documentclass[a4paper, serif, 10pt]{moderncv}

%% moderncv
% style and color
\moderncvstyle{banking}
\moderncvcolor{black}

%% page layout/margins
\usepackage[top=2.5cm, bottom=2.5cm, left=1.5cm, right=1.5cm]{geometry}

\nopagenumbers{}

%% Babel config for Brazilian Portuguese language
% ref:
% - https://latex3.github.io/babel/guides/locale-portuguese.html
\usepackage[brazilian]{babel}

%% fontspec
\usepackage{fontspec}

\IfFontExistsTF{Linux Libertine}{
  \setmainfont[Ligatures={TeX, Common}, Numbers=OldStyle]{Linux Libertine}
}{}
\IfFontExistsTF{Linux Libertine O}{
  \setmainfont[Ligatures={TeX, Common}, Numbers=OldStyle]{Linux Libertine O}
}{}

%% CTeX
% CJK support, used to show CJK characters in the resume
%
% - fontset=none: disable builtin fontset but instead set the CJK font manually
% - heading=false: disable ctex heading
% - punct=kaiming: use kaiming punctuations styles for CJK
% - scheme=plain: use plain scheme, do not override \normalsize font size
% - space=auto: space settings for CJK characters
%
% ref:
% - http://ctan.mirrorcatalogs.com/language/chinese/ctex/ctex.pdf
\usepackage[UTF8, heading=false, punct=kaiming, scheme=plain, space=auto]{ctex}

\IfFontExistsTF{Noto Serif CJK SC}{
  \setCJKmainfont{Noto Serif CJK SC}
}{}
\IfFontExistsTF{Noto Sans CJK SC}{
  \setCJKsansfont{Noto Sans CJK SC}
}{}

%% line spacing
\usepackage{setspace}
\setstretch{1.125}

%% URL styling - use normal text instead of monospace
\urlstyle{same}

% Auto-underline all links
% Defer until \begin{document} so hyperref is already loaded
\AtBeginDocument{%
  \let\oldhref\href
  \renewcommand{\href}[2]{\oldhref{#1}{\underline{#2}}}%
}

%% Patch \httplink and \httpslink to handle full URLs
% The original moderncv macros always prepend http:// or https:// to URLs.
% This causes issues when the URL already has a protocol (e.g., https://example.com)
% resulting in malformed URLs like https://https://example.com.
% This patch checks if the URL already contains "://" and uses it directly if so.
% Note: We use \str_set:Nx (x-expansion) to expand macros like \@homepage first.
\ExplSyntaxOn
\str_new:N \g_yamlresume_url_str

\RenewDocumentCommand{\httpslink}{O{}m}{
  \str_set:Nx \g_yamlresume_url_str {#2}
  \str_if_in:NnTF \g_yamlresume_url_str {://}
    { \url{#2} }
    { \url{https://#2} }
}

\RenewDocumentCommand{\httplink}{O{}m}{
  \str_set:Nx \g_yamlresume_url_str {#2}
  \str_if_in:NnTF \g_yamlresume_url_str {://}
    { \url{#2} }
    { \url{http://#2} }
}
\ExplSyntaxOff

\name{Mariana Costa}{}
\title{UX Designer Plena | Pesquisa, Prototipação e Design de Interfaces}
\phone[mobile]{+55 11 98765-4321}
\email{mariana.design@email.com}
\homepage{https://marianacosta.design}

\address{Av. Paulista, 1000, São Paulo, São Paulo, Brasil, 01310-100}{}{}

\extrainfo{{\small \faLinkedin}\ \href{https://www.linkedin.com/in/marianacosta-ux}{@marianacosta-ux} {} {} {} • {} {} {}
{\small \faBehance}\ \href{https://www.behance.net/marianacosta}{@marianacosta} {} {} {} • {} {} {}
{\small \faMedium}\ \href{https://medium.com/@marianacosta}{@marianacosta}}

\begin{document}

\maketitle

\section{Informações Básicas}

\cvline{}{UX Designer com 6+ anos de experiência em produtos digitais, atuando em todo o ciclo de discovery e delivery.
%
Especialista em pesquisa com usuários, arquitetura da informação, prototipação em alta fidelidade e design de interfaces para web e mobile. Facilita a criação e a evolução de design systems, colabora com times de produto e engenharia e conduz testes de usabilidade para embasar decisões com dados e evidências.}
\section{Formação}

\cventry{fev. de 2013–dez. de 2017}
        {Bacharelado, Design, Nota: 8.2}
        {Universidade de São Paulo (USP)}
        {\url{https://www.usp.br/}}
        {}
        {\begin{itemize}
\item Desenvolvimento de projetos práticos em design de produto e design digital
\item Pesquisa e desenvolvimento de interfaces centradas no usuário
\item Trabalhos em equipe com ênfase em metodologias ágeis e design thinking
\end{itemize}
\textbf{Cursos}: Ergonomia e Usabilidade, Design de Informação, Metodologia de Projeto em Design, Psicologia da Percepção, Tipografia, Design de Interfaces Digitais, Antropologia do Consumo, Processos de Criação em Mídias Digitais}
\section{Experiência Profissional}

\cventry{mar. de 2023–Atual}
        {UX Designer Plena}
        {Nexo Digital}
        {\url{https://nexodigital.com.br}}
        {}
        {\begin{itemize}
\item Lidera o discovery de novas funcionalidades para a plataforma de investimentos, conduzindo entrevistas, testes de usabilidade e sessões de cocriação com stakeholders
\item Cria protótipos em alta fidelidade e colabora com product managers e desenvolvedores na evolução contínua do design system
\item Implementou um processo de pesquisa contínua que reduziu em 35\% o retrabalho em especificações de produto
\item Atua como referência em acessibilidade, garantindo conformidade com WCAG 2.1 nos fluxos principais do produto
\end{itemize}
\textbf{Palavras-chave}: Design System, Pesquisa com usuários, Testes de usabilidade, Acessibilidade, Descoberta de produto, Colaboração}

\cventry{jan. de 2021–fev. de 2023}
        {UX Designer}
        {Grupo Vai}
        {\url{https://grupovai.com.br}}
        {}
        {\begin{itemize}
\item Responsável pelo redesign do aplicativo de mobilidade urbana, aumentando a taxa de conversão de cadastro em 22\%
\item Conduziu mais de 40 entrevistas e testes de usabilidade com usuários de diferentes perfis
\item Desenvolveu o design system do app, documentando padrões visuais e componentes para as plataformas Android, iOS e web
\item Trabalhou em squads multidisciplinares, participando de rituais ágeis e revisões dos componentes de interface
\end{itemize}
\textbf{Palavras-chave}: Figma, Design System, Pesquisa contextual, Prototipação, Experimentação}

\cventry{jan. de 2018–dez. de 2020}
        {Designer de Produto}
        {Up Design Studio}
        {\url{https://updesign.example}}
        {}
        {\begin{itemize}
\item Realizou projetos de branding e UI para startups dos setores de saúde, educação e fintech
\item Criou wireframes, fluxos de navegação e protótipos funcionais para validação com clientes
\item Aplicou métodos de design thinking em workshops de definição de produto
\item Colaborou com desenvolvedores na implementação dos layouts, garantindo consistência visual
\end{itemize}
\textbf{Palavras-chave}: UI Design, Wireframes, Design Thinking, Branding, Front-end}
\section{Idiomas}

\cvline{Português}{Nativo ou Bilíngue \hfill \textbf{Palavras-chave}: Português do Brasil}
\cvline{Inglês}{Proficiência Profissional Fluente \hfill \textbf{Palavras-chave}: TOEFL 102, Leitura e escrita técnicas}
\cvline{Espanhol}{Proficiência Profissional Limitada \hfill \textbf{Palavras-chave}: Comunicação básica, Leitura técnica}
\section{Competências}

\cvline{Pesquisa com Usuários}{Avançado \hfill \textbf{Palavras-chave}: Entrevistas, Testes de usabilidade, Personas, Jornada do usuário, Métricas de UX, Desk research}
\cvline{Design de Interfaces}{Especialista \hfill \textbf{Palavras-chave}: Figma, Design System, Design responsivo, Acessibilidade, Grid e espaçamento, Motion Design}
\cvline{Prototipação}{Avançado \hfill \textbf{Palavras-chave}: Alta fidelidade, Interactive prototyping, Fluxos de navegação, Design tokens}
\cvline{Colaboração em Produto}{Avançado \hfill \textbf{Palavras-chave}: Scrum, Kanban, Design critique, Documentação, Product discovery}
\cvline{Front-end para Designers}{Intermediário \hfill \textbf{Palavras-chave}: HTML, CSS, JavaScript, Tailwind CSS}
\section{Prêmios}

\cventry{nov. de 2022}
        {Associação Brasileira de Design}
        {Prêmio Latino-Americano de Design de Interfaces}
        {}
        {}
        {Reconhecimento pelo redesign do app de mobilidade urbana desenvolvido no Grupo Vai, com ênfase em usabilidade e inclusão.}

\cventry{out. de 2019}
        {Up Design Studio}
        {Destaque em Inovação}
        {}
        {}
        {Premiação interna pelo projeto de interface para uma plataforma de telemedicina, avaliado por sua navegação intuitiva e impacto positivo no tempo de atendimento.}
\section{Certificados}

\cventry{ago. de 2021}
        {Google}
        {Google UX Design Professional Certificate}
        {\url{https://www.coursera.org/professional-certificates/google-ux-design}}
        {}
        {}

\cventry{mar. de 2022}
        {Human Factors International}
        {Certified Usability Analyst}
        {\url{https://www.humanfactors.com/}}
        {}
        {}
\section{Publicações}

\cventry{jun. de 2023}
        {Case: os bastidores do redesign do fluxo de pagamento}
        {Medium}
        {\url{https://medium.com/@marianacosta/case-redesign-fluxo-pagamento}}
        {}
        {Artigo com o processo de discovery, testes de usabilidade e os resultados obtidos após a reformulação do fluxo de pagamento de um aplicativo de serviços financeiros.}

\cventry{set. de 2020}
        {Padrões de acessibilidade para aplicações móveis}
        {UXDicas}
        {\url{https://uxdicas.com.br/artigos/padroes-de-acessibilidade}}
        {}
        {Guia prático com boas práticas de acessibilidade aplicadas a aplicações móveis, incluindo contraste, leitores de tela e alvos de toque.}
\section{Referências}

\cventry{\emaillink[fernanda.lima@nexodigital.com.br]{fernanda.lima@nexodigital.com.br}}
        {Gerente de Produto na Nexo Digital}
        {Dra. Fernanda Lima}
        {+55 11 91234-1080}
        {}
        {Mariana é uma profissional extremamente empática e analítica. Sua capacidade de transformar dados de pesquisa em decisões de design elevou a qualidade do nosso produto.}
\section{Projetos}

\cventry{mar. de 2023–Atual}
        {Design system para produtos digitais da Nexo Digital}
        {Nebula Design System}
        {\url{https://marianacosta.design/projetos/nebula}}
        {}
        {Desenvolveu tokens de design, componentes acessíveis e documentação interativa para acelerar a criação de novas telas.
\textbf{Palavras-chave}: Design System, Design Tokens, Documentação, Acessibilidade}

\cventry{ago. de 2022–jan. de 2023}
        {Aplicativo de educação financeira para investidores iniciantes}
        {Jornada Invest}
        {\url{https://www.behance.net/gallery/jornada-invest}}
        {}
        {Projeto autoral de UX com pesquisa com usuários, ideação e prototipação de uma jornada gamificada para aprendizado em investimentos.
\textbf{Palavras-chave}: UX Research, Gamificação, Prototipação}
\section{Interesses}

\cvline{Acessibilidade Digital}{Tecnologia assistiva, Inclusão, WCAG}
\cvline{Fotografia}{Composição, Fotografia urbana, Edição de imagens}
\cvline{Economia Comportamental}{Heurísticas, Decisão de compra, Experiência do usuário}
\section{Voluntariado}

\cventry{jan. de 2022–Atual}
        {Mentora de Carreira}
        {Mulheres de UX Brasil}
        {\url{https://www.mulheresdeux.com.br/}}
        {}
        {Orienta mulheres em transição para a área de UX, revisa portfólios e oferece suporte em processos seletivos.}


\cventry{fev. de 2019–dez. de 2021}
        {Facilitadora de Oficinas de Design}
        {ONG Rede Incluir}
        {\url{https://www.redecincluir.org.br/}}
        {}
        {Ministrou oficinas de design participativo com jovens de escolas públicas, estimulando a criação de soluções para problemas locais.}

\end{document}